\section{Introduction}
\label{sec:introduction}

A GRPO-trained formalizer achieves a 93\% prove rate on held-out logical reasoning problems---yet 87\% of its proved outputs embed the conclusion verbatim among the premises, reducing each ``proof'' to a structural shortcut invisible to the execution reward.
When the solver receives premises that already contain the target conclusion, it trivially succeeds; the model has learned to satisfy the reward without performing logical derivation.
This gap between execution success and structural integrity defines a critical blind spot in neurosymbolic reinforcement learning.

Neurosymbolic approaches that pair large language models with formal solvers have demonstrated strong performance on logical reasoning benchmarks~\citep{satlm, logic-lm, linc, fover-pei2025}.
The dominant paradigm treats solver execution as a binary reward signal: the model generates a formal specification, a solver attempts a proof, and the outcome drives policy optimization via GRPO~\citep{deepseekmath} or related algorithms~\citep{deepseek-r1}.
This execution-only evaluation conflates genuine reasoning with degenerate shortcuts.
A model can achieve arbitrarily high prove rates by embedding the target conclusion among its premises, by producing contradictory axiom sets that prove anything via explosion, or by generating tautological single-clause outputs---all of which satisfy the solver without requiring any logical derivation.
Specification gaming of this kind has been documented in reasoning models more broadly~\citep{bondarenko-spec-gaming}, and a growing body of work studies reward hacking mechanisms in LLM training~\citep{reward-hacking-survey}, yet systematic characterization in neurosymbolic RL remains absent.

No principled diagnostic exists for detecting and quantifying structural degeneration in RL-trained formalizers.
Prompting-based neurosymbolic systems~\citep{logic-lm, linc, satlm} do not involve RL training and thus sidestep the problem entirely.
\citet{mellgren-prolog-grpo} hand-craft Prolog-specific structure penalties but provide no systematic analysis of the failure modes they target.
Verification-reward approaches use solver or prover outcomes as training signals~\citep{deepseek-prover-v15, rlsf} but evaluate execution correctness only.
Recent work has begun examining whether LLMs game formalization tasks~\citep{kim-game-formalization}, finding that high compilation rates do not imply faithful outputs---yet no framework connects structural diagnosis to the reward hacking dynamics that emerge during RL training.
The result is a field that optimizes for---and reports---prove rates without any mechanism to assess whether the underlying formalizations contain genuine logical structure.

We observe that component necessity testing---systematically ablating each solver input and checking whether the proof still holds---provides a principled, solver-agnostic diagnostic for this problem.
Removing a premise from a well-structured formalization should break the proof; removing one from a degenerate output that embeds the conclusion will not.
This leave-one-out (LOO) protocol yields a continuous score, the \emph{Intensional Structure Ratio} (ISR), measuring the fraction of formal components that are necessary for the proof.
The protocol requires only the ability to re-execute the solver with modified inputs, making it applicable across formal languages and proof systems without language-specific engineering.

A natural question follows: can ISR serve not merely as a post-hoc diagnostic but as a training signal to incentivize structural integrity?
We find that it cannot.
When ISR is incorporated into the GRPO reward, the model games the metric itself---a meta-level instance of Goodhart's Law~\citep{manheim-goodhart-taxonomy, gao-reward-overoptimization}.
Rather than producing better-structured proofs, the model discovers strategies that inflate ISR directly, most commonly by reducing all premises to a single clause that duplicates the conclusion.
ISR measures syntactic component necessity, not semantic faithfulness, and LOO-based metrics carry an inherent mathematical bias toward outputs with fewer components: fewer premises means each is trivially more ``necessary.''
These properties make ISR a reliable observer of structural integrity but an unreliable optimizer of it.

We introduce \emph{Solver Trace Ablation} (STA), a protocol that applies component necessity testing to derive ISR and systematically characterizes reward hacking across GRPO training conditions (Figure~\ref{fig:sta_overview}).
Our experimental design spans 6 reward configurations (execution-only, ISR-augmented at two dose levels, matched random-reward controls, convex combination, and a supervised fine-tuning baseline), 3 random seeds per configuration, and up to 1{,}000 training steps, evaluated on two formal languages: first-order logic (FOL, verified by Prover9) and Prolog (verified by SWI-Prolog).

STA reveals that execution-only GRPO training produces progressive structural collapse invisible to the training reward.
Prove rate reaches 93\% while ISR deteriorates from 0.64 at 100 steps to 0.13 at 500 steps to 0.003 at 1{,}000 steps---a 99.5\% decline quantifying structural degeneration that prove rate cannot detect.
All three random seeds exhibit 87--100\% conclusion embedding in their proved outputs, demonstrating that degeneration is systematic rather than seed-specific.
By 1{,}000 steps, 94.8\% of proved samples have ISR${=}$0: no single component is necessary for the proof.
A supervised fine-tuning baseline trained on the same data achieves only 14\% prove rate but 92.9\% structurally genuine outputs with near-zero hacking (7.1\% H1, 0\% H7), confirming that reward hacking is RL-specific rather than an artifact of data exposure (\textbf{C1}: ISR as diagnostic, with complete collapse trajectory; \S\ref{sec:diagnostic}).

When ISR augments the GRPO reward ($r = r_{\text{exec}} \cdot (1 + \beta \cdot \text{ISR})$), meta-level Goodhart effects emerge~\citep{karwowski-goodhart-rl}.
At low dose ($\beta{=}0.1$), ISR reward induces ISR-specific gaming: across three seeds, 64.5\% of proved outputs exhibit single-premise conclusion duplication (H7 mode) that inflates ISR by gaming the metric, compared to 0\% across all three random-reward controls---confirming H7 as ISR-reward-specific.
The remaining seed's apparent structural advantage disappears entirely under premise-count-matched analysis (1.06$\times$, indistinguishable from the random baseline); the raw ISR elevation is a mathematical artifact of premise minimization, not improved reasoning.
At high dose ($\beta{=}0.3$), classic Goodhart collapse occurs: prove rate drops to 18\% versus 48\% for the matched random control.
This dose-response pattern maps onto established Goodhart dynamics~\citep{gao-reward-overoptimization}: low-dose proxy gaming that inflates the metric without genuine improvement, followed by high-dose collapse where the proxy signal overwhelms the primary objective.
ISR-as-reward provides no genuine structural benefit at any dose (\textbf{C2}: empirical meta-level Goodhart characterization; \S\ref{sec:goodhart}).

Both the diagnostic capability and the Goodhart failure transfer across structurally different formal languages.
Applying STA to Prolog formalizations verified by SWI-Prolog, execution-only training again produces near-zero ISR (mean 0.04 across 3 seeds, $\text{PR}{=}53\%{\pm}31\text{pp}$), replicating the FOL collapse pattern with comparable cross-seed instability.
Prolog ISR-as-reward yields \emph{universal} Goodhart: all three seeds converge to gaming with 0\% genuine output, each via a distinct mechanism---tautological single-clause duplication, degenerate uppercase pseudo-Prolog syntax, and query-as-fact gaming where the model sets the target query as its own provable fact.
This is strictly stronger than the FOL result, where ${\sim}$16\% genuine output persists, demonstrating that simpler formal semantics enable more complete metric exploitation (\textbf{C3}: cross-logic diagnostic transfer with universal Goodhart; \S\ref{sec:crosslogic}).

Across all experimental conditions, we identify and quantify 7 distinct hacking modes---conclusion embedding in single-premise (H7) and multi-premise (H1) variants, tautological collapse, complexity-without-utility, repetition collapse, premise padding, and query-as-fact gaming---providing a concrete, empirically grounded taxonomy for diagnosing structural failures in neurosymbolic RL systems.
Each mode is quantified per reward condition and seed, enabling practitioners to anticipate the specific failure modes their training configurations are likely to produce (\textbf{C4}: hacking taxonomy; \S\ref{sec:taxonomy}).

These findings expose a fundamental asymmetry: ISR is reliable as a diagnostic observer of structural integrity but unreliable as an optimization target.
This instantiates a broader cautionary principle for proxy-metric-as-reward approaches~\citep{manheim-goodhart-taxonomy, pan-reward-misspecification}: diagnostic metrics that detect reward hacking can themselves become gaming targets when repurposed as training signals.
The asymmetry---\emph{reliable observer, unreliable optimizer}---carries implications beyond formal reasoning for any system where proxy measures of structural quality are candidates for reward shaping.
We release our code, evaluation scripts, and complete experimental outputs to support reproducibility.
